%=====================================================================
\clearpage
\hdr{J1}{测三棱镜顶角：反射法}

\subhdr{① 这一环交出什么}

\begin{center}
\begin{minipage}{\linewidth}
\begin{tabular}{@{}p{3.0cm}p{4.4cm}p{6.6cm}@{}}
\toprule
\thead{产出} & \thead{合格判据} & \thead{交给谁} \\
\midrule
顶角 $\alpha$ 及其不确定度 & 等边棱镜应 $59^{\circ}55'$--$60^{\circ}05'$ & J2 的折射率公式要用 $\alpha$ \\
\bottomrule
\end{tabular}
\end{minipage}
\end{center}

用时约 25 分钟：摆棱镜与调平 10 分钟，三个位置读数 15 分钟。

\subhdr{② 为什么这样能测：两束反射光的夹角是 $2\alpha$}

把一束方向固定的平行光同时照到三棱镜的两个折射面上（让棱镜的顶点对着准直管），两个面各反射一束光。\textbf{这两束反射光的夹角等于 $2\alpha$}。

\textbf{证明}：单个平面镜转过 $\gamma$，反射光线转过 $2\gamma$——这是反射定律的直接推论：入射方向不变时 $k'=k-2(k\cdot n)n$，法线转 $\gamma$，$k'$ 就转 $2\gamma$。现在两个面不是同一面镜子，但它们对同一束入射光的取向差恰好是顶角 $\alpha$，所以两束反射光的取向差是 $2\alpha$。

\warnbox{\textbf{这里有个很容易答错的地方：为什么不是 $2(180^{\circ}-\alpha)$？}棱镜两面的\textbf{外法线}夹角确实是 $180^{\circ}-\alpha$（$\alpha=60^{\circ}$ 时法线夹角 $120^{\circ}$）。但两束反射光的方向差是 $2\times120^{\circ}=240^{\circ}$，而\textbf{两条直线的夹角取小于 $180^{\circ}$ 的那个}，$360^{\circ}-240^{\circ}=120^{\circ}=2\alpha$。所以无论从哪条路算，望远镜要转的都是 $2\alpha=120^{\circ}$。}

\textbf{公式里的 $1/4$ 是两个 $\tfrac12$ 叠起来的}，不必死记：

\begin{enumerate}[label=\textbf{第 \arabic* 个 $\tfrac12$}]
\item \textbf{几何的 $\tfrac12$}：望远镜在位置 I、II 之间转过的角等于 $2\alpha$，所以 $\alpha=\tfrac12\times$ 转角；
\item \textbf{读数的 $\tfrac12$}：转角本身要用两个游标平均，$\text{转角}=\tfrac12(|{\phi_2-\phi_1}|+|{\phi_2'-\phi_1'}|)$。
\end{enumerate}

两式相乘即 $\tfrac14$。看到 $\tfrac14$ 就想「双面镜 $\times$ 双游标」。

\subhdr{③ 手怎么动}

\textbf{第 1 步：摆棱镜。}按讲义 p11 图 12-11 的方位放：让 $AB$、$AC$、$BC$ 分别垂直于三个调平螺钉两两的连线 $a_1a_2$、$a_2a_3$、$a_3a_1$；顶点稍微靠近载物台中心，这样两个面都能被准直管的光照到。

\begin{figure}[H]
\centering
\includegraphics[width=0.40\linewidth]{fig-3.12/fig12-11-prism-stage.png}
\caption{三棱镜在载物台上的方位（讲义 p11 图 12-11）：$AB\perp a_1a_2$、$AC\perp a_2a_3$、$BC\perp a_3a_1$。这样 $a_1$ 只改变 $AB$ 面法线、$a_3$ 只改变 $AC$ 面法线，两个面可以各自独立调平。}
\end{figure}

\textbf{第 2 步：调平棱镜。}用自准直法：先找到 $AB$ 面的绿十字像，调 $a_1$ 让它与上叉丝重合；再转载物台找 $AC$ 面的像，调 $a_3$ 让它重合；反复直到两面都重合。此时\textbf{棱镜主截面垂直主轴}。

\warnbox{\textbf{为什么偏偏不能调 $a_2$}：在这个摆法下 $a_1$ 只改 $AB$ 面法线、$a_3$ 只改 $AC$ 面法线，而 $a_2$ 会\textbf{同时}改动两个面——刚调好的一个面立刻被破坏，两个独立的问题被耦合成一个解不开的。讲义 p11 只用一句「Do not adjust the screw a2.」带过，原因在这里。}

\textbf{第 3 步：位置 I 读数。}把望远镜转到能看见 $AB$ 面反射的狭缝像，用\textbf{竖直叉丝}对准像的中心，锁 17 号，读\textbf{两个游标}得 $\phi_1$ 与 $\phi_1'$。

\textbf{第 4 步：位置 II 读数。}把望远镜转到 $AC$ 面的反射像，同样对准，读 $\phi_2$ 与 $\phi_2'$。一组就是这四个数。

\textbf{第 5 步：重复三组。}每组之间可以重新对一次像，但\textbf{不要动 16 号}。

\textbf{找像的先后}：如果某一面的像死活找不到，先看棱镜是否放得太偏（顶点没朝准直管）、台面高矮是否挡住光路；再看该面的调平螺钉是否把它调出了视场。

\subhdr{④ 当场记这些数}

\begin{center}
\begin{minipage}{\linewidth}
\small
\begin{tabular}{@{}p{1.4cm}p{2.7cm}p{2.7cm}p{2.7cm}p{2.7cm}p{2.0cm}@{}}
\toprule
\thead{次数} & \thead{$\phi_1$（$AB$ 面）} & \thead{$\phi_1'$} & \thead{$\phi_2$（$AC$ 面）} & \thead{$\phi_2'$} & \thead{备注} \\
\midrule
\rowstrut 1 & \blankcell & \blankcell & \blankcell & \blankcell & \\
\rowstrut 2 & \blankcell & \blankcell & \blankcell & \blankcell & \\
\rowstrut 3 & \blankcell & \blankcell & \blankcell & \blankcell & \\
\bottomrule
\end{tabular}
\end{minipage}
\end{center}

\textbf{栏目含义}：$\phi_1$ 与 $\phi_1'$ 是\textbf{同一次}望远镜位置的两个游标（差约 $180^{\circ}$），不是两次测量。若中间动过 16 号，四个数就不再对应同一个度盘，\textbf{整组作废}。

\subhdr{⑤ 当场算}

\begin{enumerate}[label=\textbf{(\arabic*)}]
\item 每行各算两个差：$|\phi_2-\phi_1|$ 与 $|\phi_2'-\phi_1'|$；\textbf{任一个超过 $180^{\circ}$，就用 $360^{\circ}$ 减它}（这条必须在每一行上分别做，不能等平均之后再判断一次）。
\item 三行的两个差各自取平均，得 $\overline{|\Delta_1|}$ 与 $\overline{|\Delta_2|}$。
\item 代公式：
\[
\alpha=\frac14\Bigl(\overline{|\phi_2-\phi_1|}+\overline{|\phi_2'-\phi_1'|}\Bigr).
\]
\end{enumerate}

\textbf{为什么是 $\tfrac14$ 而不是 $\tfrac12$}：望远镜从位置 I 转到位置 II 的角度等于 $2\alpha$，这是唯一的几何事实。两个游标各给这个转角一个估计 $|\Delta_1|$ 与 $|\Delta_2|$，取平均得
\[
\text{转角}=\frac{|\Delta_1|+|\Delta_2|}{2}=2\alpha
\quad\Longrightarrow\quad
\alpha=\frac{|\Delta_1|+|\Delta_2|}{4},
\]
即手册 (12-4)。一个 $\tfrac12$ 来自「转角是顶角的两倍」，另一个 $\tfrac12$ 来自「两个游标取平均」。

（手册把「平均」写在对四个读数上；没有折回时两者结果相同。有折回时以「每行先折回」为准，因为折回是个非线性操作。）

\textbf{例算（示例数据，非实测）}：$\phi_1=100^{\circ}12'$、$\phi_1'=280^{\circ}13'$、$\phi_2=220^{\circ}13'$、$\phi_2'=40^{\circ}14'$。

\[
|\phi_2-\phi_1|=119^{\circ}61'=120^{\circ}01',\qquad
|\phi_2'-\phi_1'|=239^{\circ}59'\xrightarrow{\ \times\ }120^{\circ}01',
\]
\[
\alpha=\frac14\bigl(120^{\circ}01'+120^{\circ}01'\bigr)=60^{\circ}00.5'.
\]

\textbf{自查三条}：① $\alpha$ 应在 $60^{\circ}$ 附近 $5'$ 以内；② 两个差的平均应约 $120^{\circ}$（$=2\alpha$）；③ 两组游标各自算出的 $\alpha$ 相差应在 $2'$ 内，差多了说明某个游标读错。

\begin{figure}[H]
\centering
\includegraphics[width=0.36\linewidth]{fig-3.12/fig12-12-vertex.png}
\caption{反射法测顶角的两个望远镜位置（讲义 p11 图 12-12）：顶点 $A$ 对住准直管，$AB$、$AC$ 两面各反射一束，望远镜在 I、II 两处读 $\phi_1,\phi_1'$ 与 $\phi_2,\phi_2'$。}
\end{figure}

\subhdr{⑥ 写进报告的一句话}

\callout{
「用反射法测得三棱镜顶角 $\alpha=60^{\circ}00.5'$（三组平均）；两反射光束夹角为 $2\alpha=120^{\circ}01'$，与等边棱镜标称值 $60^{\circ}$ 相符，偏差主要来自度盘估读与棱镜主截面未严格垂直主轴。」
}
\warnbox{
\textbf{台前只看这一段}：棱镜按图摆（$AB\perp a_1a_2$、$AC\perp a_2a_3$），只调 $a_1$、$a_3$，\textbf{不碰 $a_2$}；两面都重合后，位置 I 读两游标、位置 II 读两游标；$\alpha=\tfrac14\bigl(|\phi_2-\phi_1|+|\phi_2'-\phi_1'|\bigr)$，任一差超 $180^{\circ}$ 先折回。
}